% Part: applied-modal-logics
% Chapter: temporal-logic
% Section: language-epistemic-logic

\documentclass[../../../include/open-logic-section]{subfiles}

\begin{document}

\olfileid{aml}{tl}{sem}

\olsection{સમયસંબંધી તર્કનું અર્થવિજ્ઞાન}

\begin{defn}
સમયસંબંધી તર્કની મૂળભૂત ભાષામાં નીચેની બાબતો હોય છે:
\begin{enumerate}
  \tagitem{prvFalse}{!!{falsity} માટેનો વિધાનાત્મક અચલાંક~$\lfalse$.}{}
  \tagitem{prvTrue}{!!{truth} માટેનો વિધાનાત્મક અચલાંક~$\ltrue$.}{}
  \item !!{propositional variable}sનો !!{denumerable}s ગણ: $\Obj
    p_0$, $\Obj p_1$, $\Obj p_2$, \dots
  \item વિધાનાત્મક સંયોજકો: \startycommalist
  \iftag{prvNot}{\ycomma $\lnot$ (નિષેધ)}{}%
  \iftag{prvAnd}{\ycomma $\land$ (સંયોજન)}{}%
  \iftag{prvOr}{\ycomma $\lor$ (વિયોજન)}{}%
  \iftag{prvIf}{\ycomma $\lif$ (!!{conditional})}{}%
  \iftag{prvIff}{\ycomma $\liff$ (!!{biconditional})}{}.
  \item ભૂતકાળના સંક્રિયકો $\Ptemp$ અને $\Htemp$.
  \item ભવિષ્યના સંક્રિયકો $\Ftemp$ અને $\Gtemp$.
\end{enumerate}
\end{defn}

આગળ જઈને બીજા પ્રકારના મોડલ સંક્રિયકો ઉમેરવાની શક્યતા વિશે ચર્ચા કરીશું.

\begin{defn}
સમયસંબંધી ભાષાનાં \emph{!!^{formula}s} નીચે મુજબ આગમનાત્મક રીતે
  વ્યાખ્યાયિત થાય છે:
\begin{enumerate}
\tagitem{prvFalse}{$\lfalse$ પરમાણુક !!{formula} છે.}{}

\tagitem{prvTrue}{$\ltrue$ પરમાણુક !!{formula} છે.}{}

\item દરેક વિધાનાત્મક ચલ $\Obj p_i$ (પરમાણુક) !!{formula} છે.

\tagitem{prvNot}{જો $!A$ એ !!a{formula} હોય, તો $\lnot !A$ પણ
  !!a{formula} છે.}{}

\tagitem{prvAnd}{જો $!A$ અને $!B$ એ !!{formula}s હોય, તો $(!A \land
  !B)$ પણ !!a{formula} છે.}{}

\tagitem{prvOr}{જો $!A$ અને $!B$ એ !!{formula}s હોય, તો $(!A \lor !B)$
  પણ !!a{formula} છે.}{}

\tagitem{prvIf}{જો $!A$ અને $!B$ એ !!{formula}s હોય, તો $(!A \lif !B)$
  પણ !!a{formula} છે.}{}

\tagitem{prvIff}{જો $!A$ અને $!B$ એ !!{formula}s હોય, તો $(!A \liff !B)$
  પણ !!a{formula} છે.}{}

\sourcecorrection{OLMD-076}{મૂળમાં આ યાદીના ભવિષ્ય-સંક્રિયકને
  $F$ લખ્યો છે; ભાષાની અગાઉની યાદી અને આગળની અર્થવ્યાખ્યા મુજબ
  તેને $\Ftemp$ કર્યો છે.}
\item જો $!A$ એ !!a{formula} હોય, તો $\Ptemp  !A$, $\Htemp  !A$, $\Ftemp !A$, $\Gtemp  !A$
  એ બધાં !!{formula}s છે.

\tagitem{limitClause}{આ સિવાય બીજું કશું !!a{formula} નથી.}{}
\end{enumerate}
\end{defn}

બીજી ઇન્ટેન્શનલ તર્કપદ્ધતિઓની જેમ સમયસંબંધી તર્કપદ્ધતિઓનું
અર્થવિજ્ઞાન પણ સંબંધાત્મક નિદર્શનો વડે આપવામાં આવે છે.

\begin{defn}
  સમયસંબંધી ભાષાનું \emph{નિદર્શ} એ ત્રિક
  $\mModel{M} = \tuple{T, \prec, V}$ છે, જ્યાં
  \begin{enumerate}
  \item $T$ સમયબિંદુઓ તરીકે અર્થઘટિત થતો અખાલી ગણ છે.
  \item $\prec$ એ~$T$ પરનો દ્વિસ્થાની સંબંધ છે.
  \item $V$ એવું વિધેય છે જે દરેક !!{propositional
    variable}~$p$ને સમયબિંદુઓનો ગણ $V(p)$ આપે છે.
  \end{enumerate}
  $t \prec t'$ હોય ત્યારે કહીએ કે $t$ એ~$t'$ની \emph{પહેલાં આવે છે}.
  $t \in V(p)$ હોય ત્યારે કહીએ કે $p$ એ~$t$એ \emph{સત્ય છે}.
\end{defn}

હમણાં આપણે પૂર્વવર્તિતાના સંબંધ~$\prec$ પર કોઈ શરત મૂકતા નથી.
એટલે હાલ આપણાં સમયસંબંધી નિદર્શનોની રચના પર કોઈ નિયંત્રણ નથી:
સમય રેખીય, શાખાવાળો, ચક્રીય કે અન્ય કોઈ પણ રચનાવાળો હોય એવા
નિદર્શનો હોઈ શકે છે.

સામાન્ય મોડલ તર્કની જેમ દરેક સમયસંબંધી નિદર્શ એ નક્કી કરે છે કે
તેના કયા બિંદુએ કયાં !!{formula}s સત્ય ગણાય. સંબંધાત્મક
અર્થવિજ્ઞાનના મૂળ ખ્યાલ માટે આપણે એ જ સંજ્ઞા વાપરીએ છીએ:
``નિદર્શ~$\mModel{M}$માં !!{formula}~$!A$ બિંદુ~$t$એ સત્ય છે.''
આ સંબંધ આગમનાત્મક રીતે વ્યાખ્યાયિત થાય છે અને બધા બિનમોડલ
સંક્રિયકો માટે સામાન્ય મોડલ તર્કના કિસ્સા જેવો જ છે.

\begin{defn}\ollabel{defn:tmodels}
  નિદર્શ~$\mModel M$માં \emph{!!a{formula}~$!A$નું~$t$એ સત્ય હોવું},
  સંજ્ઞામાં $\mSat{M}{!A}[t]$, નીચે મુજબ આગમનાત્મક રીતે
  વ્યાખ્યાયિત થાય છે:
  \begin{enumerate}
  \tagitem{prvFalse}{\indcase{!A}{\lfalse}{$\mSat{M}{\lfalse}[t]$ ક્યારેય નહીં}.}{}
  \tagitem{prvTrue}{\indcase{!A}{\ltrue}{$\mSat{M}{\ltrue}[t]$ હંમેશા}.}{}
  \item $\mSat{M}{p}[t]$ ત્યારે અને માત્ર ત્યારે જ જ્યારે $t \in V(p)$
  \tagitem{prvNot}{\indcase{!A}{\lnot !B}{$\mSat{M}{\indfrm}[t]$ ત્યારે અને માત્ર ત્યારે જ જ્યારે
    $\mSat/{M}{!B}[t]$}.}{}
  \tagitem{prvAnd}{\indcase{!A}{(!B \land !C)}{$\mSat{M}{\indfrm}[t]$ ત્યારે અને માત્ર ત્યારે જ જ્યારે
    $\mSat{M}{!B}[t]$ અને $\mSat{M}{!C}[t]$}.}{}
  \tagitem{prvOr}{\indcase{!A}{(!B \lor !C)}{$\mSat{M}{\indfrm}[t]$ ત્યારે અને માત્ર ત્યારે જ જ્યારે
    $\mSat{M}{!B}[t]$ અથવા $\mSat{M}{!C}[t]$} (અથવા બંને).}{}
  \tagitem{prvIf}{\indcase{!A}{(!B \lif !C)}{$\mSat{M}{\indfrm}[t]$ ત્યારે અને માત્ર ત્યારે જ જ્યારે
    $\mSat/{M}{!B}[t]$ અથવા $\mSat{M}{!C}[t]$}.}{}
  \tagitem{prvIff}{\indcase{!A}{(!B \liff !C)}{$\mSat{M}{\indfrm}[t]$ ત્યારે અને માત્ર ત્યારે જ જ્યારે
    $\mSat{M}{!B}[t]$ અને $\mSat{M}{!C}[t]$ બંને હોય, અથવા
    $\mSat{M}{!B}[t]$ તથા $\mSat{M}{!C}[t]$ બેમાંથી એકેય ન હોય}.}{}
  \item\ollabel{defn:sub:mmodels-p}
    \indcase{!A}{\Ptemp  !B}{$\mSat{M}{\indfrm}[t]$ ત્યારે અને માત્ર ત્યારે જ જ્યારે
    $t' \prec t$ ધરાવતા કોઈ $t' \in T$ માટે $\mSat{M}{!B}[t']$}
\item\ollabel{defn:sub:mmodels-h}
    \indcase{!A}{\Htemp  !B}{$\mSat{M}{\indfrm}[t]$ ત્યારે અને માત્ર ત્યારે જ જ્યારે
    $t' \prec t$ ધરાવતા દરેક $t' \in T$ માટે $\mSat{M}{!B}[t']$}
   \item\ollabel{defn:sub:mmodels-f}
    \indcase{!A}{\Ftemp  !B}{$\mSat{M}{\indfrm}[t]$ ત્યારે અને માત્ર ત્યારે જ જ્યારે
    $t \prec t'$ ધરાવતા કોઈ $t' \in T$ માટે $\mSat{M}{!B}[t']$}
\item\ollabel{defn:sub:mmodels-g}
    \indcase{!A}{\Gtemp  !B}{$\mSat{M}{\indfrm}[t]$ ત્યારે અને માત્ર ત્યારે જ જ્યારે
    $t \prec t'$ ધરાવતા દરેક $t' \in T$ માટે $\mSat{M}{!B}[t']$}
  \end{enumerate} 
\end{defn}

આ અર્થવિજ્ઞાન પરથી જોઈ શકાય છે કે સંક્રિયકો~$\Ptemp$ અને~$\Htemp$
પરસ્પર દ્વૈત છે, અને સંક્રિયકો~$\Ftemp$ તથા~$\Gtemp$ પણ એમ જ છે.
એટલે $\Htemp  !A$ને $\lnot \Ptemp \lnot !A$ તરીકે વ્યાખ્યાયિત
કરી શકીએ છીએ; $\Gtemp$ અને~$\Ftemp$ માટે પણ આવું જ થાય છે.

\end{document}
